\documentclass[a4paper, 12pt]{article}
\usepackage[utf8]{inputenc}
\usepackage[spanish]{babel} % Cambia a 'english' si prefieres inglés
\usepackage{graphicx}
\usepackage{float}
\usepackage{amsmath}
\usepackage{booktabs}
\usepackage{array}
\usepackage{caption}
\usepackage{subcaption}
\usepackage{circuitikz} % Para diagramas de circuitos eléctricos
\usetikzlibrary{babel} % Soluciona conflictos entre circuitikz y babel-spanish
\usepackage{hyperref}
\usepackage{geometry}
\usepackage{listings}
\usepackage{xcolor}
\usepackage{siunitx} % Paquete para unidades del SI
\usepackage{enumitem} % Para continuar listas entre archivos
\usepackage{multirow} % Para tablas con celdas de varias filas
\usepackage{comment}

% Configuración de márgenes
\geometry{
    a4paper,
    left=2.5cm,
    right=2.5cm,
    top=2.5cm,
    bottom=2.5cm
}

% Configuración de hyperref
\hypersetup{
    colorlinks=true,
    linkcolor=blue,
    filecolor=magenta,      
    urlcolor=cyan,
    pdftitle={Documento Técnico de Circuitos},
    pdfpagemode=FullScreen,
}

% Configuración para listados de código (opcional)
\lstset{
    language=Python,
    basicstyle=\ttfamily\small,
    keywordstyle=\color{blue},
    commentstyle=\color{gray},
    numbers=left,
    numberstyle=\tiny\color{gray},
    frame=single,
    breaklines=true,                % Permite que las líneas largas se corten
    captionpos=b,                   % Posición del caption (abajo)
    literate=
      {á}{{\'a}}1 {é}{{\'e}}1 {í}{{\'i}}1 {ó}{{\'o}}1 {ú}{{\'u}}1
      {Á}{{\'A}}1 {É}{{\'E}}1 {Í}{{\'I}}1 {Ó}{{\'O}}1 {Ú}{{\'U}}1
      {ñ}{{\~n}}1 {Ñ}{{\~N}}1
      {¿}{{?`}}1 {¡}{{!`}}1
}

% Título del documento
\title{\textbf{TEOREMA DE NORTON}}
\author{Prof. César Augusto Rodríguez}
\date{\today}

\begin{document}

\begin{comment}
\begin{abstract}
Este documento presenta un análisis detallado del Teorema de Norton, 
aplicando sus principios a la resolución de redes de corriente continua 
con una y dos fuentes de tensión. Se incluyen desarrollos analíticos y 
ejemplos prácticos con esquemas de circuitos.
\end{abstract}
\end{comment}

\maketitle
\tableofcontents
\newpage

% --- Capítulos del Documento ---
%\section{PAGINA 1}
\section{Objetivos}
\begin{enumerate}
    \item Determinar analíticamente los valores \( I_N \) (corriente del generador de Norton) y \( R_N \) (resistencia del generador de Norton) en un circuito de c.c. que contiene una o dos fuentes de tensión.
    \item Verificar experimentalmente los valores \( I_N \) y \( R_N \) propuestos por el teorema de Norton en la solución de redes complicadas de c.c. que contienen dos fuentes de tensión.
\end{enumerate}

\section{Información Preliminar}
\section*{Teorema de Norton}
El teorema de Thévenin simplifica el análisis de las redes complicadas reduciendo el circuito original a un circuito sencillo equivalente que comprende una fuente de tensión constante, el generador de Thévenin (\( V_{TH} \)) en serie con una resistencia interna \( R_{TH} \). Este generador entrega corriente a la resistencia del circuito de carga que deseamos estudiar. Entonces es fácil determinar la corriente en la carga y la tensión aplicada entre los terminales de ésta. El teorema de Norton utiliza una técnica análoga de simplificación. Sin embargo, el generador de Norton es una fuente de corriente constante.

El teorema de Norton enuncia que cualquier red de dos terminales puede ser reemplazada por un circuito sencillo equivalente consistente en un generador de corriente constante \( I_N \), shuntado por una resistencia interna \( R_N \). La figura 27-1 a representa la red original como una unidad (bloque) terminada por una resistencia de carga \( R_L \). La figura 27-1 b muestra el circuito equivalente de Norton. La corriente de Norton \( I_N \) se distribuye entre la resistencia shunt \( R_N \) y la carga \( R_L \). La corriente \( I_L \) en \( R_L \) se puede hallar por la fórmula

\[
I_L = \frac{I_N \times R_N}{R_N + R_L}
\tag{27-1}
\]

Las reglas para determinar las constantes del circuito equivalente de Norton son las siguientes:

\begin{enumerate}
    \item La corriente constante \( I_N \) es la que debería circular en el cortocircuito entre los terminales de resistencia de carga si esta resistencia fuese reemplazada por un cortocircuito.
    \item La resistencia de Norton \( R_N \) es la que aparece entre los terminales de la carga abierta, hacia la red original, cuando las fuentes de tensión existentes en el circuito son reemplazadas por su resistencia interna. Así \( R_N \) se define de la misma manera exactamente que \( R_{TH} \) en el teorema de Thévenin.
\end{enumerate}

\section{Aplicaciones}
Consideremos el circuito de la figura 27-2 a. Deseamos hallar la corriente \( I_L \) en \( R_L \). Podemos calcular \( I_L \) mediante la aplicación de las leyes de Ohm y Kirchhoff o bien por el teorema de Thévenin, pero vamos a hallar primero \( I_L \) por el teorema de Norton.

El desarrollo del circuito equivalente de Norton para la figura 27-2 a se puede
%\section{PAGINA 2}
\begin{figure}[htbp]
\centering
\begin{subfigure}{0.8\textwidth}
    \begin{circuitikz}[scale=0.8]
        \draw (0,0) to[battery1, l={$V = 100\text{V}$}, invert] (0,3) to[R, l={$R_1 = 5\,\Omega$}] (3,3) to[R, l={$R_2 = 195\,\Omega$}] (6,3);
        \draw (6,3) -- (8,3) node[above] {F}
        to[R, l={$R_3 = 200\,\Omega$}] (8,0) node[below] {G} -- (2,0); % R3
        \draw (8,3) -- (11,3) to[R, l={$R_L = 350\,\Omega$}] (11,0) -- (8,0); % RL
        \draw (2,0) -- (0,0);
        \fill (8,3) circle (1.5pt);
        \fill (8,0) circle (1.5pt);
    \end{circuitikz}
    \caption{Circuito original}
    \label{fig:27-2-a}
\end{subfigure}

\vspace{0.5cm}
\centering
\begin{subfigure}{0.8\textwidth}
    \centering
    \begin{circuitikz}[scale=0.8]
        \draw (0,0) to[battery1, l={$V = 100\text{V}$}, invert] (0,3) to[R, l={$R_1 = 5\,\Omega$}] (3,3) to[R, l={$R_2 = 195\,\Omega$}] (6,3);
        \draw (6,3) -- (8,3) node[above] {F};
        \draw (8,3) to[ammeter, l_=$I_N$] (8,0) node[below] {G}; % I_N
        \draw (8,0) -- (0,0);
        \fill (8,3) circle (1.5pt);
        \fill (8,0) circle (1.5pt);
        \node[right, align=left] at (9, 1.5) {$I_N = \frac{100}{5+195}\text{A}$};
    \end{circuitikz}
    \caption{Cálculo de $I_N$}
    \label{fig:27-2-b}
\end{subfigure}

\vspace{0.5cm} % Add vertical space
\centering
\begin{subfigure}{0.8\textwidth}
    \centering
    \begin{circuitikz}[scale=0.8]
        \draw (2,4) -- (6,4) node[right] {F};
        \draw (2,0) -- (6,0) node[right] {G};
        \draw (2,4) to[R, l={$R_1=5\,\Omega$}] (2,2) to[R, l={$R_2=195\,\Omega$}] (2,0);
        \draw (5,4) to[R, l={$R_3=200\,\Omega$}] (5,0);
        \fill (6,4) circle (1.5pt);
        \fill (6,0) circle (1.5pt);
        \node[right, align=left] at (7.5, 2) {$R_N = \frac{200 \times (195+5)}{200+195+5}\Omega$};
    \end{circuitikz}
    \caption{Cálculo de $R_N$}
    \label{fig:27-2-c}
\end{subfigure}

\vspace{0.5cm} % Add vertical space
\centering
\begin{subfigure}{0.8\textwidth}
    \centering
    \begin{circuitikz}[scale=0.8]
        % Nodes
        \coordinate (A) at (0,3);
        \coordinate (B) at (0,0);
        % Components
        \draw (B) to[isource, l_={$I_N=0.5\text{A}$}, i>^] (A);
        \draw (A) to[short, i=$I_N$] (4,3) node[above] {F};
        \draw (4,3) to[R, l={$R_L=350\,\Omega$}] (4,0) node[below] {G} -- (B);
        \fill (4,3) circle (1.5pt);
        \fill (4,0) circle (1.5pt);
    \end{circuitikz}
    \caption{Circuito Equivalente de Norton con Carga}
    \label{fig:27-2-d}
\end{subfigure}

\caption{Aplicación del teorema de Norton al Análisis de una Red de c.c.}
\label{fig:27-2}
\end{figure}

\noindent seguir fácilmente en las figuras 27-2b, c y d. En la figura 27-2b, $R_L$ ha sido cortocircuitada. En consecuencia $R_3$ también lo está. La corriente de Norton $I_N$ en el cortocircuito será pues:
\[
I_N = \frac{V}{R_T} = \frac{100}{195+5}\text{A} = 0,5\text{A}
\]
En la figura 27-2c vemos que la fuente original de tensión V ha sido reemplazada por su resistencia interna $R_1$ que es igual a 5$\Omega$. La resistencia $R_N$ es la que aparece desde los terminales de la carga abierta, hacia la red original cuando las fuentes de tensión existentes en el circuito son reemplazadas por su resistencia interna. $R_N$ es $100\Omega$. Finalmente, en la figura 27-2d se ve que el generador equivalente de Norton $I_N$ está shuntado con $R_N$, en paralelo con la resistencia de carga $R_L$. Aplicando la fórmula (27-1) tenemos:
\[
I_L = \frac{I_N \times R_N}{R_N + R_L} = \frac{0,5 \times 100}{100 + 350}\text{A} = 0,111\text{A}
\]
Este es el mismo valor obtenido en la práctica 26 cuando hemos hecho uso del Teorema de Thevenin para hallar $I_L$.

Ahora será posible hallar la corriente correspondiente a cualquier valor de $R_L$ en el circuito de la figura 27-2a por medio del generador equivalente de Norton con valores de la figura 27-2d.

\section{Resolución de las  Redes de c.c. con DOS o MAS Generadores.}
La resolución de redes de c.c. complicadas que contienen dos o más fuentes de tensión es posible utilizando las leyes de Kirchhoff, el Teorema de Thévenin y el Teorema de Norton. La aplicación de los dos últimos teoremas simplifica a menudo el análisis y los cálculos. A modo de ejemplo vamos a resolver el circuito de la figura 27-3 por los tres métodos.

\begin{figure}[htbp]
\centering
\begin{subfigure}{0.8\textwidth}
    \centering
    \begin{circuitikz}[scale=0.8]
        \draw (0,3) node[anchor=east] {A} to[battery1, l={$V_S$}] (0,0) node[anchor=east] {B};
        \draw (3,3) to[short, i=$I$] (0,3); \draw (3,3) to[R, l={$R_2$}] (6,3);
        \draw (6,3) -- (8,3) node[above] {F};
        \draw (0,0) to[short, i=$I$] (3,0) to[R, l={$R_1$}] (6,0) -- (8,0) node[below] {G};
        \draw (8,3) -- (11,3) to[battery1, l={$V_T$}] (11,0) -- (8,0); % VT
        \fill (8,3) circle (1.5pt);
        \fill (8,0) circle (1.5pt);
        \fill (0,3) circle (1.5pt);
        \fill (0,0) circle (1.5pt);
    \end{circuitikz}
    \caption{Circuito original}
    \label{fig:27-3-a}
\end{subfigure}
\caption{Aplicación de la Ley de Tensión de Kirchhoff a un Circuito Cerrado}
\label{fig:27-3}
\end{figure}
%\section{PAGINA 3}

Por las leyes de Kirchhoff. Recordemos las dos leyes, que son:

\begin{enumerate}
    \item La suma de las caídas de tensión en resistencias conectadas en serie en un circuito cerrado es igual a la tensión aplicada. Se puede enunciar ahora la ley en otra forma, a saber, la suma algébrica de las tensiones en cualquier bucle o circuito cerrado es nula.
    
    \item La corriente que entra en cualquier nudo de un circuito eléctrico es igual a la corriente que sale del nudo. También se puede enunciar esta ley en otra forma, a saber, la suma algébrica de las corrientes en cualquier punto de un circuito es nula.
\end{enumerate}

En la aplicación de la segunda forma de la ley de tensión de Kirchhoff hay que atenerse a determinados convenios. Consideremos el circuito de la figura 27-3. La expresión «suma algébrica», etc., requiere poner una cuidadosa atención en los signos de las fuentes de tensión y de las caídas de tensión. Así, empezando en A en la figura 27-3 y recorriendo el bucle ABCDA, tenemos primero la tensión \( V_s \). Esta es positiva en A. Por tanto, el primer término de la ecuación es \(+ V_s\). Luego tenemos la caída de tensión \( V_1 \) entre los extremos de \( R_1 \). Suponiendo que el sentido del flujo de electrones sea el indicado por la flecha, la polaridad de la tensión \( V_1 \) es la indicada en la figura 27-3, o sea de - a +. De aquí que el término siguiente de la ecuación sea \(- V_1\). Procediendo análogamente en el sentido CDA, las polaridades de las tensiones son, respectivamente, \(- V_T\) y \(- V_2\). La ecuación que satisface la ley de tensión de Kirchhoff se puede escribir ahora como sigue:
\[
+ V_s - V_1 - V_T - V_2 = 0 \tag{27-5}
\]

\begin{center}
\textit{NOTA: Los signos opuestos de \( V_s \) y \( V_T \) indican que estas baterías están conectadas opuestamente, o sea que sus tensiones no se suman, hecho que resulta evidente por el esquema del circuito.}
\end{center}

Apliquemos las leyes de Kirchhoff al circuito de la figura 27-4. Es necesario expresar tantas ecuaciones como caminos independientes de corriente existen (bucles). El bucle 1 es ABCDA. Suponemos que el flujo de electrones \( I_1 \) tiene en este bucle el sentido indicado por la flecha. Otro bucle completo independiente, el 2, es DFGCD y se supone que el flujo de electrones \( I_2 \) es el indicado. \( I_1 \) e \( I_2 \) son las

\begin{figure}[htbp]
    \centering
    \begin{circuitikz}[american, thick, scale=0.8]
        \draw (0,0) node[below] {B} to[battery1, l={$V_1=\SI{100}{\volt}$}, invert] (0,4) node[above] {A};
        \draw (0,4) to[R, l={$R_1=\SI{1}{\kilo\ohm}$}] (5,4) node[above] {D};
        \draw (5,4) to[battery1, l={$V_2=\SI{150}{\volt}$}] (5,2) to[R, l={$R_2=\SI{1.2}{\kilo\ohm}$}] (5,0) node[below] {C};
        \draw (5,0) -- (0,0);
        \draw (5,4) -- (10,4) node[above] {F} to[R, l={$R_3=\SI{1.2}{\kilo\ohm}$}] (10,0) node[below] {G} -- (5,0);
        \fill (5,4) circle (1.5pt);
        \fill (5,0) circle (1.5pt);
        \fill (0,4) circle (1.5pt);
        \fill (0,0) circle (1.5pt);
        \fill (10,4) circle (1.5pt);
        \fill (10,0) circle (1.5pt);
        \draw[red, thick, ->] (2.5,2) ++(-220:0.7) arc (-220:120:0.7);
        \node[red] at (2.5,2) {$I_1$};
        \draw[blue, thick, ->] (8,2) ++(120:0.7) arc (120:-220:0.7);
        \node[blue] at (8,2) {$I_2$};
    \end{circuitikz}
    \caption{Circuito Original}
    \label{fig:27-4a}
\end{figure}

\begin{figure}[htbp]
    \centering
    \begin{circuitikz}[american, thick, scale=0.8]
        \draw (1,0) node[left] {B} to[battery1, l={$V_1=\SI{100}{\volt}$}, invert] (1,4) node[left] {A};
        \draw (1,4) -- (2,4) to[R, l={$R_1=\SI{1}{\kilo\ohm}$}] (5,4) -- (7,4) node[above] {D};
        \draw (7,4) to[battery1, l={$V_2=\SI{150}{\volt}$}, invert] (7,2) to[R, l={$R_2=\SI{2.2}{\kilo\ohm}$}] (7,0);
        \draw (7,4) -- (9.7,4) node[above] {F}; \draw (7,0) -- (9.7,0) node[below] {G};
        \draw (7,0) node[below] {C} -- (1,0);
        \fill (9.7,4) circle (1.5pt);
        \fill (9.7,0) circle (1.5pt);
        \fill (1,4) circle (1.5pt);
        \fill (1,0) circle (1.5pt);
        \fill (7,4) circle (1.5pt);
        \fill (7,0) circle (1.5pt);
        \draw[green, thick, ->] (4,2) ++(-220:0.7) arc (-220:120:0.7);
        \node[green] at (4,2) {$I_3$};
    \end{circuitikz}
    \caption{Circuito para el Cálculo de la Tensión de Thévenin.}
    \label{fig:27-4b}  
\end{figure}

\begin{figure}[htbp]
    \centering
    \begin{circuitikz}[american, thick, scale=0.8]
        \draw (1,0) node[below] {B} to[short] (1,4) node[above] {A} -- (3,4) to[R, l={$R_1=\SI{1.0}{\kilo\ohm}$}] (7,4);
        \draw (7,4) -- (9,4) node[above] {D} -- (11,4) node[above] {F};
        \draw (9,4) to[R, l={$R_2=\SI{2.2}{\kilo\ohm}$}] (9,0) node[below] {C} -- (11,0) node[below] {G};
        \draw (9,0) -- (1,0); % Mantengo la línea original a B
        \fill (1,0) circle (1.5pt);
        \fill (1,4) circle (1.5pt);
        \fill (11,4) circle (1.5pt);
        \fill (11,0) circle (1.5pt);
        \fill (9,4) circle (1.5pt);
        \fill (9,0) circle (1.5pt);
    \end{circuitikz}
    \caption{Circuito para el Cálculo de la Resistencia Equivalente.}
    \label{fig:27-4c}
\end{figure}

\begin{figure}[htbp]
    \centering
    \begin{circuitikz}[american, thick, scale=0.8]
        \draw (1,0) to[battery1, l={$V_{TH}=\SI{21.82}{\volt}$}, invert] (1,3) to[R, l={$R_{TH}=\SI{687.5}{\ohm}$}] (1,6) -- (8,6);
        \draw (8,6) -- (9,6) node[above] {F} to[R, l={$R_3=\SI{1.2}{\kilo\ohm}$}] (9,0) node[below] {G};
        \draw (9,0) -- (1,0);
        \fill (9,6) circle (1.5pt);
        \fill (9,0) circle (1.5pt);
    \end{circuitikz}
    \caption{Circuito Equivalente de Thévenin con la Carga $R_3$}
    \label{fig:27-4d}
\end{figure}

\noindent En la Figura \ref{fig:27-4a} se muestra el circuito de dos mallas con los bucles ABCDA y DFGCD. Se han indicado las corrientes de lazo \(I_1\) (en rojo) e \(I_2\) (en azul), ambas en sentido horario. La Figura \ref{fig:27-4b} presenta el equivalente de Thévenin que se obtendría al simplificar la red desde un par de terminales.
%\section{PAGINA 4}
\noindent incógnitas de este sistema de ecuaciones. Finalmente se supone que se desea hallar 
$I_2$ en $R_3$. Empezamos en A en el lazo 1 y escribimos la ecuación de tensión: 
\[
+ V_1 - (I_1 + I_2) \times R_2 + V_2 - I_1 \times R_1 = 0
\tag{27-6}
\]
\noindent Sustituyendo los valores de circuito en la ecuación (27-6), tenemos:
\[
100 - 2.200 \times (I_1 + I_2) + 150 - 1.000 \times I_1 = 0
\tag{27-7}
\]
\noindent Simplificando, obtenemos:
\[
3.200I_1 + 2.200I_2 = 250
\tag{27-8}
\]

\noindent Ahora empezamos en el punto D del bucle 2. La ecuación de tensión es:
\[
- I_2R_3 - (I_1 + I_2)R_2 + 150 = 0
\tag{27-9}
\]

\noindent Substituyendo los valores del circuito en la ecuación (27-9):
\[
- 1.200I_2 - 2.200(I_1 + I_2) + 150 = 0
\tag{27-10}
\]

\noindent Simplificando, tenemos:
\[
2.200I_1 + 3.400I_2 = 150
\tag{27-11}
\]

\noindent Con el sistema de dos ecuaciones (27-8) y (27-11) hallamos $I_1$ e $I_2$, obteniendo:
\[
\begin{aligned}
I_1 &= 0,0861 \, \text{A} \\
I_2 &= -0,0115 \, \text{A}
\end{aligned}
\]  

\noindent El signo negativo en $I_2$ significa simplemente que la corriente tiene sentido contrario al supuesto inicialmente. $I_2$ es la corriente que circula por $R_3$.

\subsection{Resolución de circuito por métodos de análisis}

\subsection{Enunciado del problema}
Consideremos el circuito de la figura 27-4 donde se desea hallar la corriente $I_2$ que circula por la resistencia $R_3$.

\subsection{Resolución por leyes de Kirchhoff}
Empezamos en el punto A del bucle 1 y escribimos la ecuación de tensión:
\[
+ V_1 - (I_1 + I_2)R_2 + V_2 - I_1 \times R_1 = 0
\tag{27-6}
\]

Substituyendo los valores del circuito en la ecuación (27-6), tenemos:
\[
100 - 2.200(I_1 + I_2) + 150 - 1.000I_1 = 0
\tag{27-7}
\]

Simplificando, obtenemos:
\[
3.200I_1 + 2.200I_2 = 250
\tag{27-8}
\]

Ahora empezamos en el punto D del bucle 2. La ecuación de tensión es:
\[
- I_2R_3 - (I_1 + I_2)R_2 + 150 = 0
\tag{27-9}
\]

Substituyendo los valores del circuito en la ecuación (27-9):
\[
- 1.200I_2 - 2.200(I_1 + I_2) + 150 = 0
\tag{27-10}
\]

Simplificando, tenemos:
\[
2.200I_1 + 3.400I_2 = 150
\tag{27-11}
\]

Con el sistema de dos ecuaciones (27-8) y (27-11) hallamos $I_1$ e $I_2$, obteniendo:
\[
\begin{aligned}
I_1 &= 0,0861 \, \text{A} \\
I_2 &= -0,0115 \, \text{A}
\end{aligned}
\]  

\noindent El signo negativo en $I_2$ significa simplemente que la corriente tiene sentido contrario al supuesto inicialmente. $I_2$ es la corriente que circula por $R_3$.

\section*{Resolución de circuito por métodos de análisis}

\subsection{Enunciado del problema}
Consideremos el circuito de la figura 27-4 donde se desea hallar la corriente $I_2$ que circula por la resistencia $R_3$.

\subsection{Resolución por leyes de Kirchhoff}
Empezamos en el punto A del bucle 1 y escribimos la ecuación de tensión:
\[
+ V_1 - (I_1 + I_2)R_2 + V_2 - I_1 \times R_1 = 0
\tag{27-6}
\]

Substituyendo los valores del circuito en la ecuación (27-6), tenemos:
\[
100 - 2.200(I_1 + I_2) + 150 - 1.000I_1 = 0
\tag{27-7}
\]

Simplificando, obtenemos:
\[
3.200I_1 + 2.200I_2 = 250
\tag{27-8}
\]

Ahora empezamos en el punto D del bucle 2. La ecuación de tensión es:
\[
- I_2R_3 - (I_1 + I_2)R_2 + 150 = 0
\tag{27-9}
\]

Substituyendo los valores del circuito en la ecuación (27-9):
\[
- 1.200I_2 - 2.200(I_1 + I_2) + 150 = 0
\tag{27-10}
\]

Simplificando, tenemos:
\[
2.200I_1 + 3.400I_2 = 150
\tag{27-11}
\]

Con el sistema de dos ecuaciones (27-8) y (27-11) hallamos $I_1$ e $I_2$, obteniendo:
\[
\begin{aligned}
I_1 &= 0,0861 \, \text{A} \\
I_2 &= -0,0115 \, \text{A}
\end{aligned}
\]

El signo negativo en $I_2$ significa simplemente que la corriente tiene sentido contrario al supuesto inicialmente. $I_2$ es la corriente que circula por $R_3$.

\subsection{Resolución por el teorema de Thévenin}
Para hallar $I_2$ en $R_3$ aplicando el teorema de Thévenin, primero calculamos $V_{TH}$, que es la tensión en circuito abierto entre los puntos F y G.

La tensión de circuito abierto entre F y G es $V_{FG} = V_{TH} = -150 + I_3R_2$. Por tanto, necesitamos hallar $I_3$ antes de determinar $V_{TH}$.

Aplicando la ley de tensión de Kirchhoff al bucle ABCDA:
\[
100 - I_3(2.200) + 150 - I_3(1.000) = 0
\tag{27-12}
\]

Resolviendo:
\[
\begin{aligned}
3.200I_3 &= 250 \\
I_3 &= \frac{250}{3.200} = 0,0781 \, \text{A}
\end{aligned}
\]

Entonces:
\[
I_3R_2 = 0,0781 \times 2.200 = 171,82 \, \text{V}
\]

La tensión Thévenin es:
\[
V_{TH} = -150 + I_3R_2 = 21,82 \, \text{V}
\tag{27-13}
\]

Para hallar $R_{TH}$, substituimos $V_1$ y $V_2$ por sus resistencias internas (supuestas iguales a cero) y hallamos la resistencia entre F y G:
\[
R_{TH} = \frac{2.200 \times 1.000}{2.200 + 1.000} = 687,5 \, \Omega
\tag{27-14}
\]

Finalmente, hallamos $I_2$ en el circuito equivalente de Thévenin:
\[
I = I_2 = \frac{V_{TH}}{R_{TH} + R_3} = \frac{21,82}{687,5 + 1.200} = \frac{21,82}{1.887,5} = 0,0115 \, \text{A}
\tag{27-15}
\]

\subsection{Resolución por el teorema de Norton}
Para hallar $I_2$ en $R_3$ por el teorema de Norton, desarrollamos el circuito equivalente como se muestra en las figuras 27-5b y 27-5c.

En la figura 27-5b, $R_3$ está reemplazada por un cortocircuito. $I_N$, la corriente de cortocircuito entre los terminales F y G, es la corriente del generador equivalente de Norton.

Para hallar $I_N$ aplicamos las leyes de Kirchhoff a los bucles 1 y 2 de la figura 27-5b:
\begin{itemize}
    \item Bucle 1: Trayecto ABCGFDA
    \item Bucle 2: Trayecto DFGCD
\end{itemize}

\subsection{Conclusiones}
\begin{itemize}
    \item Se obtiene el mismo resultado ($I_2 = 0,0115 \, \text{A}$) por los tres métodos
    \item El teorema de Thévenin y Norton proporcionan métodos sistemáticos para simplificar circuitos complejos
    \item La elección del método depende de la configuración específica del circuito y de lo que se necesite calcular
\end{itemize}

\begin{table}[H]
    \centering
    \caption{Comparación de resultados obtenidos por diferentes métodos}
    \label{tab:comparacion_metodos}
    \begin{tabular}{lcc}
        \toprule
        \textbf{Método} & \textbf{Corriente $I_2$ (A)} & \textbf{Observaciones} \\
        \midrule
        Leyes de Kirchhoff & 0,0115 & Método general pero laborioso \\
        Teorema de Thévenin & 0,0115 & Útil para análisis de un elemento \\
        Teorema de Norton & 0,0115 & Equivalente a Thévenin, diferente formulación \\
        \bottomrule
    \end{tabular}
\end{table}

\textbf{Nota:} Los valores numéricos corresponden a resistencias en ohmios, tensiones en voltios y corrientes en amperios.
%\section{PAGINA 5}

\section{Análisis Usando el Teorema de Norton}

A continuación, se presentan los circuitos que utilizaremos:

\begin{figure}[htbp]
    \centering
    \begin{circuitikz}[european]
        \draw (0,0) node[below] {B} to[battery, l={$V_s=\SI{100}{\volt}$}, invert] (0,4) node[above] {A} to[R, l={$R_1=\SI{1}{\kilo\ohm}$}] (4,4) node[above] {D};
        \draw (4,4) to[battery, l={$V_2=\SI{150}{\volt}$}] (4,2) to[R, l={$R_2=\SI{2.2}{\kilo\ohm}$}] (4,0) node[below] {C};
        \draw (4,4) -- (8,4) node[above] {F};
        \draw (8,4) to[R, l={$R_3=\SI{1.2}{\kilo\ohm}$}] (8,2) to[ammeter, l=$I_2$] (8,0) node[below] {G};
        \draw (8,0) -- (4,0) -- (0,0); % Conecta G -> C -> B
        \fill (0,0) circle (1.5pt); % Nodo B
        \fill (0,4) circle (1.5pt); % Nodo A
        \fill (4,4) circle (1.5pt); % Nodo D
        \fill (4,0) circle (1.5pt); % Nodo C
        \fill (8,4) circle (1.5pt); % Nodo F
        \fill (8,0) circle (1.5pt); % Nodo G
    \end{circuitikz}
    \caption{Solución de una Red Compleja por el Teorema de Norton.}
    \label{fig:circuito original}
\end{figure}

\begin{figure}[htbp]
    \centering
    \begin{circuitikz}[european]
        \draw (0,0) node[below] {B} to[battery, l={$V_s=\SI{100}{\volt}$}, invert] (0,4) node[above] {A} to[R, l={$R_1=\SI{1}{\kilo\ohm}$}] (4,4) node[above] {D};
        \draw (4,4) to[battery, l={$V_2=\SI{150}{\volt}$}] (4,2) to[R, l={$R_2=\SI{2.2}{\kilo\ohm}$}] (4,0) node[below] {C};
        \draw (4,4) -- (8,4) node[above] {F};
        \draw (8,4) to[ammeter, l=$I_N$] (8,0) node[below] {G}; % R3 y I2 reemplazados por I_N
        \draw (8,0) -- (4,0) -- (0,0); % Conecta G -> C -> B
        \fill (0,0) circle (1.5pt); % Nodo B
        \fill (0,4) circle (1.5pt); % Nodo A
        \fill (4,4) circle (1.5pt); % Nodo D
        \fill (4,0) circle (1.5pt); % Nodo C
        \fill (8,4) circle (1.5pt); % Nodo F
        \fill (8,0) circle (1.5pt); % Nodo G
    \end{circuitikz}
    \caption{Cálculo de la corriente de Norton ($I_N$) cortocircuitando los terminales F y G.}
    \label{fig:circuito-in}
\end{figure}

\begin{figure}[htbp]
    \centering
    \begin{circuitikz}[european]
        \draw (0,0) node[below] {B} to[short] (0,4) node[above] {A} to[R, l={$R_1=\SI{1}{\kilo\ohm}$}] (4,4) node[above] {D};
        \draw (4,4) to[R, l={$R_2=\SI{2.2}{\kilo\ohm}$}] (4,0) node[below] {C};
        \draw (4,4) -- (8,4) node[above] {F};
        \draw (8,0) -- (4,0) -- (0,0); % Conecta G -> C -> B
        \fill (0,0) circle (1.5pt); % Nodo B
        \fill (0,4) circle (1.5pt); % Nodo A
        \fill (4,4) circle (1.5pt); % Nodo D
        \fill (4,0) circle (1.5pt); % Nodo C
        \fill (8,4) circle (1.5pt); % Nodo F
        \node at (8,0) [below] {G}; % Nodo G como terminal abierto
    \end{circuitikz}
    \caption{Cálculo de la resistencia de Norton ($R_N$) vista desde los terminales F y G.}
    \label{fig:circuito-rn}
\end{figure}

\begin{figure}[htbp]
    \centering
    \begin{circuitikz}[european]
        \draw (0,0) node[below] {B} to[isource, l={$I_N=\SI{31.8}{\milli\ampere}$}] (0,4) node[above] {A} to[short] (4,4) node[above] {D};
        \draw (4,4) to[R, l={$R_N=\SI{687.5}{\ohm}$}] (4,0) node[below] {C};
        \draw (4,4) -- (8,4) node[above] {F};
        \draw (8,4) to[R, l={$R_3=\SI{1.2}{\kilo\ohm}$}] (8,2) to[ammeter, l=$I_N$] (8,0) node[below] {G};
        \draw (8,0) -- (4,0) -- (0,0); % Conecta G -> C -> B
        \fill (0,0) circle (1.5pt); % Nodo B
        \fill (0,4) circle (1.5pt); % Nodo A
        \fill (4,4) circle (1.5pt); % Nodo D
        \fill (4,0) circle (1.5pt); % Nodo C
        \fill (8,4) circle (1.5pt); % Nodo F
        \fill (8,0) circle (1.5pt); % Nodo G
    \end{circuitikz}
    \caption{Circuito Equivalente de Norton con la carga conectada.}
    \label{fig:circuito-equivalente-norton}
\end{figure}

\noindent En el lazo 1,

\begin{equation}
\begin{split}
    &100 - I_4 \times (1000) = 0 \\
    &I_4 = \frac{100}{1000} = \SI{0.1}{\ampere}
\end{split} \tag{27-16}
\end{equation}

En el bucle 2, 

\begin{equation}
\begin{split}
    &2200 \times I_5 + 150 = 0 \\
    &I_5 = \frac{150}{2200} = \SI{0.0682}{\ampere}
\end{split} \tag{27-17}
\end{equation}

La figura 27-5b muestra que $I_N=I_4-I_5$, puesto que los sentidos de 
$I_4$ y $I_5$ son opuestos. Por lo tanto,

\begin{equation}
\begin{split}
    &I_N = I_4 - I_5 = \SI{0.1}{\ampere} - \SI{0.0682}{\ampere} = \SI{0.0318}{\ampere}
\end{split} \tag{27-18}
\end{equation}  

La resistencia del generador de Norton $R_N$ es la resistencia total de 
la red original, es decir, la resistencia total de la red original es 
$R_N=\SI{687.5}{\ohm}$. Por lo tanto, el generador de Norton 
es $R_N=\SI{687.5}{\ohm}$ y $I_N=\SI{0.0318}{\ampere}$.

Nuevamente hallamos que el valor de la corriente en $R_3$ es el mismo que 
el valor obtenido utilizando las leyes de Kirchhoff y el teorema de 
Thévenin.

\section{Resumen}
\begin{itemize}[series=resumenlist]
\item El teorema de Norton ofrece otro método de resolver circuitos complicados. 
Por el teorema de Norton se puede reemplazar una red complicada de dos 
terminales por un circuito equivalente sencillo que actúa como el circuito 
original en la carga conectada a los dos terminales.
\item El circuito equivalente se compone de un generador de corriente, llamado
generador Norton ($I_N$) en paralelo con la resistencia interna $R_N$ del 
generador de Norton y en paralelo con los cuales está también conectada 
la carga $R_L$. La corriente $I_N$ del generador se divide entonces entre 
$R_N$ y $R_L$. La figura 27-1b muestra el generador de Norton y su carga. 
Es el equivalente de Norton de un circuito complicado tal como el representado 
en forma de bloques en la figura 27-1a.
\item Para determinar la corriente de Norton $I_N$, se cortocircuita la carga 
y se calcula la corriente que fluye por el circuito. Esta corriente de cortocircuito 
es $I_N$. Para calcular $I_N$ puede ser necesario hacer uso de las leyes de Ohm 
y de Kirchhoff.
\item Para determinar la resistencia de Norton $R_N$ que actúa en paraleo con 
el generador de Norton se aplica la misma técnica que para hallar la 
la resistencia de Thévenin en el experimento precedente. Se procedde como 
sigue: Abrir la carga en los dos terminales en cuestión de la red original. Reemplazar 
todas la fuentes de tensión por sus resistencias internas. Luego calcular $R_N$, 
resistencia en los terminales abiertos de la carga, hacia el circuito.
\item El teorema de Norton se aplica a un circuito con una o mas fuentes de 
alimentación.
\item Una vez reemplazada la red complicada por el circuito equivalente de Norton 
se puede hallar la corriente $I_L$ que pasa por la carga aplicando la fórmula 
siguiente (fig. 27-1b):

\begin{equation}
\begin{split}
    &I_L = \frac{I_N \times R_N}{R_N + R_L}
\end{split} \tag{}
\end{equation}
\end{itemize}
%\section{PAGINA 6}

%\section{Continuación del Resumen}
\begin{itemize}[series=resumenlist, resume]
    \item Este es un nuevo ítem que continúa la lista de la página anterior.
    \item Y aquí puedes agregar otro ítem más.
\end{itemize}

\section{Autoexamen}
\begin{itemize}
    \item Considerar el circuito de la figura 27-2a. $V=120 V$, $R_1=10$, 
    $R_2=390 \Omega$, $R_3=600 \Omega$, $R_L=270 \Omega$. Se supone que 
    la resistencia interna de la batería V es 0 en el circuito 
    equivalente de Norton.
    \begin{enumerate}
        \item $I_N = ..................A$.
        \item $R_N = .................. \Omega$.
        \item $I_L = ..................A$.
    \end{enumerate}
    \item Se considera el circuito de la figura 27-5a. $V_1=45 V$, $V_2=75 V$, y 
    se supone que la resistencia interna de estas baterías de cero, 
    $R_1=450 \Omega$, $R_2=1500 \Omega$, $R_3 (la carga) =470 \Omega$.
    En el circuito equivalente de Norton,
    \begin{enumerate}
        \item $I_N = ..................A$.
        \item $R_N = .................. \Omega$.
        \item $I_L = ...................A$.
    \end{enumerate}
\end{itemize}

\section{Materiales Necesarios}
\begin{itemize}
    \item Equipo: Dos fuentes de c.c. de tensión regulada.
    \item Dos miliamperímetros de varios alcances o dos VOM capaces de 
    medir hasta \SI{50}{\milli\ampere}.
    \item EVM.
    \item Resistencias: \SI{0.5}{\watt} \SI{3300}{\ohm}, \SI{0.5}{\watt} 
    \SI{5100}{\ohm}, \SI{10000}{\ohm} y otras que requieran la operación 8 (discrecional).
    \item Diversos: Potenciómetros 10000 $\Omega$ 2 W.
\end{itemize}

\section{Procedimiento}
\begin{enumerate}[series=procedimientolist]
    \item Conectar el circuito de la figura 27-6a. $S_1$ está abierto. 
    Conectar las fuentes de alimentación reguladas $V_1$ y $V_2$ y 
    ajustar $V_1=40 V$ y $V_2=20 V$. 
    \item $S_2$ está abierto. Cerrar $S_1$. Cerciorarse de que $V_1$ y 
    $V_2$ se matienen en $40 V$ y $20 V$ respectivamente. Medir con el 
    amperímetro M la corriente $I_L$ en $R_3$. Anotarla en la tabla 27-1, en 
    la columna \textbf{$I_L$ medida circuito original}.
    \item Cerrar $S_2$ cortocircuitando $R_3$. Medir la corriente $I_N$ que 
    suministra la fuente $V_1$ al cortocircuito. Anotarla en la tabla 27-1, en la 
    columna \textbf{$I_N$ medida circuito original}.
    \item Suprimir $V_1$ y $V_2$ y reemplazarlas por sus resistencias internas 
    (cortocircuitos). Eliminar $R_3$ en el circuito. Cortocicuitar los nodos 
    D y F entre si, así como los A y B. $S_1$ está aun cerrado, $S_2$ abierto. 
    Medir la resistencia entre los nodos F y G y anotarla en la tabla 27-1, en la 
    columna \textbf{$R_N$ medida circuito original}. Esta es la resistencia en shunt 
    con el generador equivalente de Norton.
    \item Conectar el circuito de la figura 27-6b. $S_1$ está abierto. 
    $R_3$ es la resistencia de carga de 5100 ohmios del circuito anterior. 
    $M_1$ y $M_2$ son miliamperímetros para medir la corriente de Norton 
    y la corriente de carga respectivamente. Ajustar la salida de tensión 
    V en cero voltios.
\end{enumerate}
\input{capitulos/TEOREMA_NORTON-07.tex}
%\section{PAGINA 8}

\begin{table}[h!]
    \centering
    \caption{Resultados del experimento.}
    \label{tab:27-1}
    \begin{tabular}{|c|c|c|c|c|c|c|c|}
        \hline
        \multirow{3}{*}{$R_L$} & \multicolumn{2}{c|}{$I_N$ (\si{\milli\ampere})} & \multicolumn{2}{c|}{$R_N$ (\si{\ohm})} & \multicolumn{3}{c|}{$I_L$ (\si{\milli\ampere})} \\
        \cline{2-8}
        & \multirow{2}{*}{Medida} & \multirow{2}{*}{Calculada} & \multirow{2}{*}{Medida} & \multirow{2}{*}{Calculada} & \multicolumn{2}{c|}{Medida} & \multirow{2}{*}{Calculada} \\
        \cline{6-7}
        & & & & & Norton & Circuito original & \\
        \hline
        \SI{5100}{\ohm} & & & & & & & \\
        \hline
         & & & & & & & \\
        \hline
         & & & & & & & \\
        \hline
         & & & & & & & \\
        \hline
    \end{tabular}
\end{table}

\vspace{2em}

\begin{enumerate}[label=\arabic*., resume]
    \setcounter{enumi}{6} % Continuamos la lista en el número 7
    \item Si se utilizan en el experimento valores de $R_3$ (fig. 27-6 a) distintos de \SI{5100}{\ohm}, ¿cuál debe ser el valor de $I_N$ en la figura 27-6 b? Explicarlo.
\end{enumerate}

\vspace{1cm}

\section{Respuestas a las preguntas de autoexamen}
\begin{enumerate}[label=\arabic*.]
    \item 
    \begin{enumerate}[label=\alph*)]
        \item 0,3
        \item 240
        \item 0,141
    \end{enumerate}
    \item 
    \begin{enumerate}[label=\alph*)]
        \item 0,05
        \item 346
        \item 0,021
    \end{enumerate}
\end{enumerate}

\end{document}